% =====================================================================
%  preamble.tex —— 导言区（Lead 独占；译者只读，不得修改）
%  植田一石《ホモロジー的ミラー対称性》（东京大学 2015 年度冬学期讲义录）
%  中文译注本 · XeLaTeX 编译
% =====================================================================
\documentclass[11pt,a4paper]{article}

\usepackage[UTF8]{ctex}
\usepackage{amsmath,amssymb,amsthm,mathtools}
\usepackage{bm}
\usepackage{graphicx}
\usepackage[a4paper,left=2.4cm,right=2.4cm,top=2.6cm,bottom=2.7cm]{geometry}
\usepackage[dvipsnames]{xcolor}
\usepackage{enumitem}
\usepackage{booktabs}
\usepackage{array}
\usepackage{multicol}
\usepackage{fancyhdr}
\usepackage{caption}
\usepackage{titlesec}
\usepackage[hidelinks,unicode]{hyperref}
\hypersetup{pdftitle={同调镜像对称（中文译注本）},
  pdfauthor={植田一石 著；中文译注：本工程}}

\captionsetup{font=small,labelsep=quad}
\setlength{\emergencystretch}{3em}
\allowdisplaybreaks[1]
\graphicspath{{figures/}{../work/figures/}{work/figures/}}

% ---------- 节标题：原文「第 N 節」 ----------
\titleformat{\section}{\Large\bfseries}{第 \thesection{} 节}{1em}{}
\titlespacing*{\section}{0pt}{22pt}{12pt}
\titleformat{\subsection}{\large\bfseries}{\thesubsection}{1em}{}
\titleformat{\subsubsection}[runin]{\bfseries}{\thesubsubsection.}{0.5em}{}

% ---------- 页眉页脚 ----------
\pagestyle{fancy}
\fancyhf{}
\fancyhead[L]{\small\itshape Homological Mirror Symmetry}
\fancyhead[R]{\small 同调镜像对称\quad\thepage}
\renewcommand{\headrulewidth}{0.4pt}

% ---------- 定理类环境：全节共享一个计数器 ----------
% 原文：定義 30.1 / 定理 30.3 / 命題 30.4 …（同一节内连续编号）
\theoremstyle{plain}
\newtheorem{theorem}{定理}[section]
\newtheorem{definition}[theorem]{定义}
\newtheorem{proposition}[theorem]{命题}
\newtheorem{lemma}[theorem]{引理}
\newtheorem{corollary}[theorem]{推论}
\newtheorem{example}[theorem]{例}
\newtheorem{claim}[theorem]{断言}
\newtheorem{conjecture}[theorem]{猜想}
\newtheorem{question}[theorem]{问题}
\newtheorem{exercise}[theorem]{习题}
\theoremstyle{definition}
\newtheorem{remark}[theorem]{注}
\newtheorem{notation}[theorem]{记号}
\newtheorem{construction}[theorem]{构造}

% ---------- 常用宏 ----------
\newcommand{\R}{\mathbb{R}}
\newcommand{\C}{\mathbb{C}}
\newcommand{\Z}{\mathbb{Z}}
\newcommand{\N}{\mathbb{N}}
\newcommand{\Q}{\mathbb{Q}}
\newcommand{\PP}{\mathbb{P}}
\newcommand{\kk}{k}
\newcommand{\dd}{\mathrm{d}}
\newcommand{\ee}{\mathrm{e}}
\newcommand{\ii}{\mathrm{i}}
\newcommand{\id}{\mathrm{id}}
\newcommand{\rank}{\operatorname{rank}}
\newcommand{\Hom}{\operatorname{Hom}}
\newcommand{\End}{\operatorname{End}}
\newcommand{\Ext}{\operatorname{Ext}}
\newcommand{\Tot}{\operatorname{Tot}}
\newcommand{\Spec}{\operatorname{Spec}}
\newcommand{\Coh}{\operatorname{Coh}}
\newcommand{\Perf}{\operatorname{Perf}}
\newcommand{\Dcoh}{D^{b}(\Coh)}
\newcommand{\Fuk}{\mathcal{F}}
\newcommand{\cA}{\mathcal{A}}\newcommand{\cB}{\mathcal{B}}\newcommand{\cC}{\mathcal{C}}
\newcommand{\cD}{\mathcal{D}}\newcommand{\cE}{\mathcal{E}}\newcommand{\cF}{\mathcal{F}}
\newcommand{\cG}{\mathcal{G}}\newcommand{\cH}{\mathcal{H}}\newcommand{\cI}{\mathcal{I}}
\newcommand{\cL}{\mathcal{L}}\newcommand{\cM}{\mathcal{M}}\newcommand{\cN}{\mathcal{N}}
\newcommand{\cO}{\mathcal{O}}\newcommand{\cP}{\mathcal{P}}\newcommand{\cS}{\mathcal{S}}
\newcommand{\cT}{\mathcal{T}}\newcommand{\cU}{\mathcal{U}}\newcommand{\cV}{\mathcal{V}}
\newcommand{\cW}{\mathcal{W}}\newcommand{\cX}{\mathcal{X}}\newcommand{\cY}{\mathcal{Y}}
\newcommand{\cZ}{\mathcal{Z}}
\newcommand{\fg}{\mathfrak{g}}\newcommand{\fM}{\mathfrak{M}}\newcommand{\fT}{\mathfrak{T}}
\newcommand{\fX}{\mathfrak{X}}\newcommand{\fY}{\mathfrak{Y}}\newcommand{\fZ}{\mathfrak{Z}}

% 原文行内粗体小标题（如「**定義.**」）用这个
\newcommand{\para}[1]{\par\noindent\textbf{#1}\quad}

% 图注：原文「図 13.1: ...」→ 印「图 13.1：...」
\newcommand{\figtitle}[2]{\caption*{图 #1：\ #2}}
% 表注
\newcommand{\tabtitle}[2]{\caption*{表 #1：\ #2}}

% 译者注框（审校批注就地引用时用；正文不滥用）
\newcommand{\transnote}[1]{\par\smallskip\noindent{\small\color{Sepia}【译者注】#1}\smallskip\par}

\setlength{\parindent}{2em}
\setcounter{tocdepth}{2}
